\section{Mechanisms nobody designed to be optimal}
\label{sec:experiment}

Everything above describes records chosen \emph{optimally}. No coordinator
chooses one that way: it keeps a count, or a threshold, or the first few
sub-agents' reports. We evaluate such mechanisms by enumerating every profile, so
the numbers are exact---no sampling, no seed, no error bar---and we compute the
margin \emph{via} Theorem~\ref{thm:elim}, one LP per agent rather than a search
over payment rules. The experiment therefore uses the paper's main tool as its
instrument. Predictions were registered and hashed before the harness was
written; the falsified ones are reported in Appendix~\ref{app:episodes}.

\begin{table}[t]
\centering
\caption{Compaction mechanisms at $n = 5$, uniform prior. ``Flat'' means
$D = \Dnull(K)$ \emph{and} $\gamma > 0$. The last two rows are the same size and
the same distortion.}
\label{tab:compact}
\small
\begin{tabular}{@{}lrrrrl@{}}
\toprule
mechanism & $K$ & $D$ & $\Dnull(K)$ & $\gamma$ & flat \\
\midrule
count of completions & $6$  & $1.5625$ & $0.8750$ & $0.2000$ & no \\
first $4$ agents only & $16$ & $0.5000$ & $0.5000$ & $0.0000$ & no \\
\midrule
flat construction & $10$ & $0.6875$ & $0.6875$ & $0.5000$ & \textbf{yes} \\
$K$-median code   & $10$ & $0.6875$ & $0.6875$ & $0.0000$ & no \\
\bottomrule
\end{tabular}
\end{table}

\paragraph{The last two rows are the paper.}
\texttt{flat} and \texttt{kmedian} use the same ten labels and incur \emph{the
same} distortion $0.6875 = \Dnull(10)$. Both are reconstruction-optimal at their
budget. One admits a payment rule with margin $1/2$; the other admits none at
all, at any price, for any bounded decoder. Compression quality is therefore not
a proxy for accountability even at fixed size: two records can be
indistinguishable on every reconstruction metric and differ by everything on the
incentive axis. A designer who selects a compaction rule by measuring
reconstruction error alone has no information about which of these two they
built.

\paragraph{Two exact values worth naming.}
An aggregate count has margin exactly $1/n$---not $0$, which is what we
registered, and not $1$. The prediction of $1$ came from the point-mass edge
condition; the minimax over \emph{mixtures} in Theorem~\ref{thm:elim} telescopes
and cancels the intermediate shifts, and $1/n$ is what survives. Fixed truncation
to the first $m$ agents has margin exactly $0$: the dropped agents are invisible,
so by Corollary~\ref{cor:edge} no payment reaches them, and the margin is a
minimum over agents. Retaining a strict subset of agents' bits, chosen in
advance, is worth exactly nothing.

\paragraph{Where the theory was checked rather than illustrated.}
No mechanism beats the frontier: every one sits on or above $\Dnull(K)$ for its
own $K$, every two-label mechanism sits on or above Theorem~\ref{thm:onebit}, and
no mechanism below $K = \gtwo(\Qn)$ is flat. These were derivations---their
failure would have meant a bug or a false theorem---but they are checked against
records nobody tuned. Appendix~\ref{app:episodes} carries a second experiment
over sampled episodes; Section~\ref{sec:llm} leaves chosen distributions behind
altogether and puts the same instruments on artifacts a language model wrote.
